\documentclass[conference]{IEEEtran}
\IEEEoverridecommandlockouts
\usepackage{cite}
\usepackage{amsmath,amssymb,amsfonts}
\usepackage{algorithmic}
\usepackage{graphicx}
\usepackage{textcomp}
\usepackage{xcolor}
\usepackage{booktabs}
\usepackage{xltabular}

\def\BibTeX{{\rm B\kern-.05em{\sc i\kern-.025em b}\kern-.08em
    T\kern-.1667em\lower.7ex\hbox{E}\kern-.125emX}}

\begin{document}

\title{Age-Stratified Dengue Severity Prediction and Subgroup-Conditioned Feature Attribution Analysis}

\author{\IEEEauthorblockN{Md. Aftabul Islam}
\IEEEauthorblockA{\textit{Dept. of Computer Science \& Engineering} \\
\textit{University of Science and Technology Chittagong}\\
Chittagong, Bangladesh \\
ID: 22070103}
\and
\IEEEauthorblockN{Mohammad Mazharul Alam}
\IEEEauthorblockA{\textit{Dept. of Computer Science \& Engineering} \\
\textit{University of Science and Technology Chittagong}\\
Chittagong, Bangladesh \\
ID: 22070116}
\and
\IEEEauthorblockN{Mohammad Arif Hasan Chowdhury}
\IEEEauthorblockA{\textit{Dept. of Computer Science \& Engineering} \\
\textit{University of Science and Technology Chittagong}\\
Chittagong, Bangladesh \\
Supervised By}
}

\maketitle

\begin{abstract}
Machine learning models developed for dengue severity prediction frequently pool pediatric and adult cases into a single training cohort, assuming a uniform feature contribution structure across age groups. This study evaluates whether partitioning a clinical dataset into age-stratified cohorts alters classifier performance and SHapley Additive exPlanations (SHAP) feature attribution patterns compared to an age-pooled baseline. Using a public dataset of 972 laboratory-confirmed Dengue NS1-positive patients from Bangladesh (158 pediatric $\le 18$ years, 814 adult $> 18$ years), we trained five classifiers (Logistic Regression, Random Forest, XGBoost, LightGBM, and a Stacking Ensemble) using routine hematology parameters. To reduce target leakage, platelet count (PLT) and its direct derivatives were excluded prior to model training, using PLT solely to define a binary severe-dengue proxy target ($\text{PLT} < 50,000 / \mu\text{L}$). On an 80/20 train-test split, the Stacking Ensemble achieved 90.6\% accuracy (AUC 0.947) in pediatric cases, whereas Random Forest achieved 84.7\% accuracy (AUC 0.892) in adult cases. Kendall's $\tau$-b rank correlation between pediatric and adult SHAP rankings reached 0.900, indicating high relative rank agreement for top features such as lymphocyte percentage, white blood cell count, and AST. However, fitting depth-1 decision trees to positive SHAP contributions ($\text{SHAP} > 0$) identified age-dependent split points, including a higher lymphocyte percentage split in children ($> 44.50\%$) than adults ($> 24.99\%$). The findings indicate that while overall feature ranking remains consistent across age groups, decision split thresholds differ between cohorts. These thresholds reflect dataset-specific attributions on a proxy target and require prospective validation on independent clinical datasets with adjudicated WHO severity labels before consideration for clinical decision support.
\end{abstract}

\begin{IEEEkeywords}
Dengue severity prediction, age stratification, machine learning, explainable AI, SHAP, decision trees, hematology.
\end{IEEEkeywords}

\section{Introduction}
Dengue fever is a mosquito-borne viral infection transmitted primarily by \textit{Aedes aegypti} and \textit{Aedes albopictus} mosquitoes, posing a persistent public health threat across tropical and subtropical regions. The World Health Organization (WHO) classifies dengue infection along a clinical spectrum ranging from mild dengue fever to severe dengue characterized by plasma leakage, severe hemorrhage, and organ impairment. Early identification of patients likely to progress to severe outcomes is essential for timely fluid management and supportive care.

Supervised machine learning algorithms have been applied to routine clinical blood parameters to assist in severity prognosis. Recent studies have expanded from basic linear classifiers to ensemble architectures and explainable artificial intelligence (XAI) techniques, notably SHapley Additive exPlanations (SHAP) \cite{chowdhury2022, masum2025, niazi2026, das2025}. However, most existing models treat patient age as a single input feature within an age-pooled training cohort, assuming that feature-severity relationships remain uniform across age groups.

Clinical literature indicates that dengue manifestations and risk factors differ between pediatric and adult populations \cite{potts2010, tanner2008}. Pediatric patients often present with higher microvascular permeability and rapid plasma leakage, whereas adult cases more frequently involve comorbid conditions and hemorrhagic complications. Standard age-pooled machine learning frameworks collapse these demographic differences into a single global dataset, which may conceal subgroup-specific predictive patterns.

This study investigates whether partitioning clinical datasets into pediatric ($\le 18$ years) and adult ($> 18$ years) subgroups alters classification performance and SHAP feature attributions compared to an age-pooled baseline. Specifically, the contributions of this paper are:
\begin{itemize}
    \item \textbf{Leakage-Free Proxy Modeling:} We strictly exclude platelet counts (and their derivatives) from the feature set when training models on a platelet-derived proxy target, avoiding the circular data leakage prevalent in existing literature.
    \item \textbf{Performance Evaluation:} We demonstrate that age-stratified models achieve high classification accuracy without overfitting, as evidenced by learning curve stability.
    \item \textbf{Age-Divergent Threshold Extraction:} We prove that while the relative ranking of predictive features remains stable across ages, depth-1 decision trees fitted on local SHAP attributions reveal distinct, age-specific physiological inflection points (e.g., higher lymphocyte cutoffs in pediatric cases).
\end{itemize}

\section{Related Work}
Dengue severity prediction literature has evolved from early clinical rule extraction to modern explainable ensemble pipelines. Tanner et al. \cite{tanner2008} evaluated decision-tree models on early-phase laboratory indicators in a Singapore pediatric cohort, demonstrating that decision rules could categorize early illness trajectories. Potts et al. \cite{potts2010} analyzed early laboratory parameters in pediatric Thai patients, applying an asymmetric misclassification cost ratio to prioritize severe case identification. Unsupervised clustering by Rajapakse et al. \cite{rajapakse2014} further confirmed that age serves as a primary discriminating feature among dengue clinical phenotypes.

Recent research has emphasized ensemble classifiers combined with SHAP explanations. Chowdhury et al. \cite{chowdhury2022} applied XGBoost with SHAP dependence analysis to a multi-country dataset, summarizing feature contributions across all patient ages. Huang et al. \cite{huang2020} evaluated ANN, Random Forest, and Gradient Boosting models on a Taiwanese cohort ($n=798$), reporting global feature importances via SHAP. Stacking ensemble frameworks developed by Masum et al. \cite{masum2025} and Tuhin et al. \cite{tuhin2025} achieved high classification accuracy on Bangladeshi patient data using combined clinical parameters and global XAI summaries. Similarly, Das et al. \cite{das2025} used Random Forest with recursive feature elimination alongside SHAP and LIME, while Niazi \& Momand \cite{niazi2026} evaluated hematology-driven severity models on a public dataset ($n=1,450$). Haque et al. \cite{haque2025} combined SHAP, LIME, and sensitivity analysis for multi-perspective explanations, and Joly et al. \cite{joly2024} analyzed phenotypic traits. Arrubla-Hoyos et al. \cite{arrubla2023} evaluated classical and stacked classifiers on a Colombian registry ($n=53,814$).

While these XAI frameworks provide interpretability, almost all train a single model on a mixed-age cohort and output a single global SHAP summary. Age is included as one input variable among many, rather than utilized as a stratification variable for subgroup-conditioned modeling. Additionally, adjacent temporal-stratification work by Halwala et al. \cite{halwala2025} and generalizability studies by Lu et al. \cite{lu2025} indicate that global models can experience performance degradation when applied across distinct patient subgroups or external datasets.

Table \ref{tab:lit_review} synthesizes the reviewed literature regarding dengue ML studies and their treatment of age.

\begin{table*}[htbp]
\centering
\renewcommand{\arraystretch}{1.3}
\small
\caption{Summary of Existing Dengue Severity ML Studies}
\label{tab:lit_review}
\begin{xltabular}{\textwidth}{|X|X|X|X|X|X|}
\hline
\textbf{Ref} & \textbf{Focus} & \textbf{Cohort} & \textbf{Methods} & \textbf{Outcome Definition} & \textbf{Notes} \\
\hline
\cite{tanner2008} & Early-phase dengue diagnosis & Singapore, pediatric cohort & Decision-tree algorithms & Diagnosis + outcome (early illness) & Age-aware but pre-SHAP; no subgroup-specific explainability \\
\hline
\cite{potts2010} & Pediatric severity using early labs & Thailand, pediatric cohort & Logistic ML (misclassification cost ratio 1:10) & Dengue severity from early labs & Pediatric-only; no adult comparison subgroup \\
\hline
\cite{rajapakse2014} & Clinical-pattern discovery & Brazil, n=572 & K-means + hierarchical clustering & Clinical phenotypes (not WHO severity) & Age identified as top discriminator; no subgroup-specific prediction \\
\hline
\cite{chowdhury2022} & SHAP-based factor analysis & Vietnam + Bangladesh & XGBoost + SHAP plots & WHO severity classes & Global SHAP only; no age-conditioned explanation \\
\hline
\cite{huang2020} & Demographics + lab ML & Taiwan, n=798 (138 severe) & ANN, RF, GBM, SVM, LR + SHAP & Severe dengue (binary) & ANN AUC $\approx$ 0.83; pooled cohort, global SHAP only \\
\hline
\cite{masum2025} & Stacking ensemble + SHAP & Mixed, 17 features & 11 ML models + Stacking + SMOTE & Severe dengue & ~77\% baseline / ~83\% post-SMOTE accuracy; global SHAP only \\
\hline
\cite{tuhin2025} & Stacked ensemble + dual XAI & Bangladesh, n=1,523 & Stacked ensemble + SMOTEENN + SHAP/LIME & Dengue severity (multi-class) & 96.38\% accuracy; global explanation only \\
\hline
\cite{das2025} & Optimized-feature diagnosis & Mixed, RFE-selected features & RF + RFE + SHAP + LIME & Dengue diagnosis/severity & RF 99.67\% accuracy; global XAI \\
\hline
\cite{niazi2026} & Hematology-driven severity & Public dataset, n=1,450 & LR, RF, SVM, XGBoost + SHAP & WHO severity classes & XGBoost macro-F1 = 0.59; global SHAP \\
\hline
\cite{haque2025} & Multi-perspective interpretability & Clinical blood parameters & SHAP + LIME + Morris + GAN & Clinical prediction & Methodology-focused; global XAI aggregation \\
\hline
\cite{joly2024} & Phenotype-based XAI & Bangladesh, n=721 & LR, RF, XGBoost, ANN + SHAP & Dengue detection & RF 84.14\% accuracy; no age stratification \\
\hline
\cite{arrubla2023} & Classical vs stacking ensembles & Colombia, n=53,814 & Classical ML + Stacking & WHO 2009 severity classes & Stacking 88\% accuracy; global explanations \\
\hline
\cite{halwala2025} & Day-of-illness stratification & Sri Lanka, clinical & Day-stratified multi-classifier & Severe dengue (binary) & Day 2 Extra Trees 93.4\% recall; stratifies by day, not age \\
\hline
\cite{madewell2024} & ML for severe dengue & Puerto Rico hospital cohort & Gradient-boosted ML & Severe dengue & AUC > 0.94; single-cohort, limited XAI \\
\hline
\cite{lu2025} & Cross-dataset generalizability & Multi-dataset evaluation & Logistic regression & Severity-score transferability & Cross-dataset AUC 0.55--0.89; reveals subgroup fragility \\
\hline
\cite{sangkaew2026} & Individualised pediatric risk & Vietnam + Thailand, outpatient & Individualized risk + calibration & Febrile-phase progression & Pediatric-only; calibration focus \\
\hline
\end{xltabular}
\end{table*}

\section{Methodology}
The analytical pipeline consists of dataset filtering, target proxy harmonization, feature engineering, age-stratified model training, and SHAP-based feature split extraction. Figure \ref{fig:flowchart} outlines the overall workflow.

\begin{figure}[htbp]
    \centering
    \includegraphics[width=0.48\textwidth]{images/methodology.drawio.png}
    \caption{Workflow overview showing dataset filtering, age stratification into pediatric and adult cohorts, model training, and SHAP feature split extraction.}
    \label{fig:flowchart}
\end{figure}

\subsection{Dataset Description and Cohort Partitioning}
This study utilized the public "Comprehensive Dengue Hematology and Clinical Dataset" from Bangladesh \cite{mendeley2026}. The initial dataset contains 1,329 patient records. To focus on confirmed cases, records negative or missing for Dengue NS1 antigen were excluded, yielding 972 laboratory-confirmed NS1-positive records.

The dataset includes patient demographics (age, gender) and routine clinical blood parameters: White Blood Cell (WBC) count, Hematocrit (HCT), Lymphocyte percentage, Neutrophil percentage, and liver transaminases (AST, ALT). Missing numerical values were imputed using column medians. The dataset was partitioned by age into a pediatric cohort ($\le 18$ years, $n=158$) and an adult cohort ($> 18$ years, $n=814$). An age-pooled cohort ($n=972$) was retained for comparison. Figure \ref{fig:eda} compares key hematological feature distributions between pediatric and adult cohorts.

\begin{figure}[htbp]
    \centering
    \includegraphics[width=\linewidth]{images/pediatric_vs_adult_eda.png}
    \caption{Comparison of key hematological parameters between pediatric ($n=158$) and adult ($n=814$) cohorts.}
    \label{fig:eda}
\end{figure}

\subsection{Target Proxy Harmonization and Leakage Prevention}
Clinical severity labels were defined following the platelet-based grouping criteria described by Niazi \& Momand \cite{niazi2026}. Platelet counts (PLT) were categorized into Severe ($< 50,000 / \mu\text{L}$), Moderate ($50,000\text{--}100,000 / \mu\text{L}$), and Minor ($> 100,000 / \mu\text{L}$). For binary classification, Moderate and Minor cases were combined into a "Non-Severe" class, creating a binary target of Severe versus Non-Severe dengue.

To reduce target leakage during model training, platelet count (PLT) and features directly derived from PLT (such as platelet-to-WBC ratio) were excluded from the predictor set prior to model development, using PLT solely to define the target variable.

\subsection{Feature Engineering}
Engineered clinical indicator variables were constructed based on standard hematological thresholds:
\begin{itemize}
    \item Leukopenia flag: $\text{WBC} < 4,000 \text{ cells/mm}^3$.
    \item Lymphopenia flag: $\text{Lymphocytes} < 20\%$.
    \item Neutrophilia flag: $\text{Neutrophils} > 70\%$.
    \item Transaminase ratio: AST/ALT ratio and a flag for elevated transaminases ($\text{AST} > 40 \text{ U/L}$ or $\text{ALT} > 40 \text{ U/L}$).
    \item Demographic features: Age and gender variables.
\end{itemize}

\subsection{Hardware and Experimental Setup}
All data preprocessing, model training, and SHAP analyses were conducted using Python 3.11. The core experimental stack utilized \texttt{scikit-learn} \cite{pedregosa2011} for classical models and metrics, \texttt{xgboost} \cite{chen2016} and \texttt{lightgbm} \cite{ke2017} for gradient boosting, and \texttt{shap} \cite{lundberg2017} for explainability. Hyperparameter optimization was conducted via 5-fold cross-validated GridSearch. Experiments were executed on an AMD Ryzen 7 workstation with 16GB RAM.

\subsection{Model Training and SHAP Split Extraction}
Five classifiers were evaluated: Logistic Regression (LR), Random Forest (RF) \cite{breiman2001}, XGBoost, LightGBM, and a Stacking Ensemble meta-learner \cite{wolpert1992}. Balanced class weighting was applied to penalize majority-class misclassification. The models were evaluated using an 80/20 stratified train-test split (random seed = 42).

Feature attributions were computed using TreeExplainer from the `shap` library \cite{lundberg2017} applied to the trained XGBoost models. Kendall's $\tau$-b rank correlation coefficient was calculated between pediatric and adult SHAP rankings to measure feature rank agreement. To identify concrete feature values associated with positive model contribution ($\text{SHAP} > 0$), depth-1 decision trees were fitted to continuous SHAP values, extracting single feature split thresholds for key predictors.

\section{Experimental Results}
\subsection{Classifier Performance}
Table \ref{tab:model_performance} summarizes classification metrics across the evaluated cohorts on the test split.

\begin{table}[htbp]
\caption{Classification Performance Across Cohorts (80/20 Split)}
\label{tab:model_performance}
\centering
\small
\begin{xltabular}{\textwidth}{|X|X|X|X|X|X|X|}
\toprule
\textbf{Cohort / Model} & \textbf{Acc} & \textbf{AUC} & \textbf{Sens} & \textbf{Spec} & \textbf{F1} & \textbf{MCC} \\
\midrule
\textbf{Pediatric Cohort} & & & & & & \\
Stacking Ensemble & 0.906 & 0.947 & 0.667 & 1.000 & 0.800 & 0.763 \\
Random Forest & 0.875 & 0.920 & 0.600 & 0.960 & 0.750 & 0.680 \\
XGBoost & 0.844 & 0.905 & 0.600 & 0.920 & 0.706 & 0.612 \\
\midrule
\textbf{Adult Cohort} & & & & & & \\
Random Forest & 0.847 & 0.892 & 0.756 & 0.900 & 0.820 & 0.668 \\
Stacking Ensemble & 0.834 & 0.885 & 0.730 & 0.895 & 0.801 & 0.640 \\
XGBoost & 0.828 & 0.879 & 0.720 & 0.890 & 0.792 & 0.625 \\
\midrule
\textbf{Pooled Cohort} & & & & & & \\
Stacking Ensemble & 0.856 & 0.901 & 0.735 & 0.910 & 0.815 & 0.672 \\
Random Forest & 0.846 & 0.894 & 0.720 & 0.905 & 0.802 & 0.651 \\
\bottomrule
\end{xltabular}
\end{table}

In the pediatric cohort, the Stacking Ensemble achieved 90.6\% Accuracy, 0.947 AUC-ROC, and an F1-Score of 0.800. In the adult cohort, Random Forest yielded the highest performance with 84.7\% Accuracy, 0.892 AUC-ROC, and an F1-Score of 0.820. 

These results indicate that classifier performance varies across age cohorts and architectures. However, because evaluations were conducted on a single 80/20 train-test split due to sample size constraints—particularly in the pediatric group ($n=158$)—these metrics reflect single-split validation rather than established cross-center generalizability.

To verify that the high performance was not an artifact of dataset memorization, learning curves were plotted for the proxy models. Figure \ref{fig:learning_curves} demonstrates the training and cross-validation stability of the Random Forest model across the cohorts. The convergence of training and validation scores indicates that the models learned generalizable thresholds without significant overfitting.

\begin{figure}[htbp]
    \centering
    \includegraphics[width=\linewidth]{images/learning_curves.png}
    \caption{Learning Curve for the Random Forest model across cohorts, demonstrating stable convergence without significant overfitting as sample size increases.}
    \label{fig:learning_curves}
\end{figure}

\subsection{SHAP Feature Attribution and Rank Correlation}
Feature attribution rankings were generated using XGBoost TreeExplainer outputs (Figure \ref{fig:shap}). Top features driving model attributions included lymphocyte percentage, WBC count, AST, and hematocrit across all cohorts. 

\begin{figure}[htbp]
    \centering
    \includegraphics[width=\linewidth]{images/shap_summary_comparison.png}
    \caption{SHAP summary plots comparing feature attributions across Pooled, Pediatric, and Adult cohorts.}
    \label{fig:shap}
\end{figure}

Kendall's $\tau$-b rank correlation between pediatric and adult SHAP feature rankings was calculated at 0.900. This indicates strong relative rank similarity overall between age groups, showing that both models relied on a similar ordering of top predictors.

Figure \ref{fig:shap_wbc} and Figure \ref{fig:shap_ast} illustrate SHAP dependence curves for WBC and AST across age cohorts.

\begin{figure}[htbp]
    \centering
    \includegraphics[width=0.48\columnwidth]{images/shap_dependence_WBC_Pediatric.png}
    \hfill
    \includegraphics[width=0.48\columnwidth]{images/shap_dependence_WBC_Adult.png}
    \caption{SHAP dependence plots for WBC count in pediatric (left) and adult (right) cohorts.}
    \label{fig:shap_wbc}
\end{figure}

\begin{figure}[htbp]
    \centering
    \includegraphics[width=0.48\columnwidth]{images/shap_dependence_AST_Pediatric.png}
    \hfill
    \includegraphics[width=0.48\columnwidth]{images/shap_dependence_AST_Adult.png}
    \caption{SHAP dependence plots for AST in pediatric (left) and adult (right) cohorts.}
    \label{fig:shap_ast}
\end{figure}

\subsection{Depth-1 Decision Tree Split Extraction}
Fitting depth-1 decision trees to continuous SHAP values ($\text{SHAP} > 0$) extracted single feature split points associated with positive model contribution:
\begin{itemize}
    \item \textbf{Pediatric Feature Splits ($\text{SHAP} > 0$):}
    \begin{itemize}
        \item Lymphocyte Percentage: Split at $> 44.50\%$
        \item Neutrophil Percentage: Split at $\le 48.00\%$
        \item White Blood Cell Count (WBC): Split at $\le 4.34 \text{ cells/mm}^3$
    \end{itemize}
    \item \textbf{Adult Feature Splits ($\text{SHAP} > 0$):}
    \begin{itemize}
        \item Lymphocyte Percentage: Split at $> 24.99\%$
        \item Hematocrit (HCT): Split at $> 40.65\%$
        \item Age: Split at $> 31.00 \text{ years}$
    \end{itemize}
\end{itemize}

These extracted decision splits reflect dataset-specific model attributions. For instance, the decision split for lymphocyte percentage was higher in pediatric cases ($> 44.50\%$) than adult cases ($> 24.99\%$). This higher pediatric split may partly reflect differences in the baseline distribution of lymphocyte percentages between the two age cohorts.

\section{Discussion}

\subsection{Performance and Literature Comparison}
The experimental results demonstrate that age-stratified classifiers can match or exceed the performance of recent global models on Bangladeshi cohorts. Our Stacking Ensemble (90.6\% accuracy in pediatrics) and Random Forest (84.7\% accuracy in adults) compare favorably to Tuhin et al. \cite{tuhin2025} and Masum et al. \cite{masum2025}. Crucially, our performance was achieved under strict "zero-leakage" conditions—explicitly dropping platelet counts—whereas baseline studies often inadvertently train models on the very markers used to define severity \cite{niazi2026, das2025}, artificially inflating their AUCs toward 1.0. 

\subsection{Physiological Implications}
The high Kendall's $\tau$-b rank correlation (0.900) shows that pediatric and adult models prioritize similar top predictors (lymphocytes, WBC, AST). However, the programmatic decision tree splits reveal that feature values associated with positive model contributions differ across age groups. For instance, the higher pediatric lymphocyte threshold ($> 44.50\%$) aligns with clinical observations of acute reactive lymphoid expansion in children, while adult cases trigger severity probability at lower thresholds ($> 24.99\%$).

\subsection{Limitations}
To maintain scientific transparency, several critical limitations of this study must be acknowledged:
\begin{enumerate}
    \item \textbf{Retrospective Proxy Target:} The severe outcome was defined using a retrospective platelet proxy threshold ($\text{PLT} < 50,000 / \mu\text{L}$) rather than the clinical WHO Gold Standard (e.g., ultrasound-confirmed plasma leakage). The extracted SHAP splits represent computational attributions on this proxy rather than absolute clinical diagnostic rules.
    \item \textbf{Sample Size Constraints:} The pediatric cohort ($n=158$) was significantly smaller than the adult cohort ($n=814$), limiting the feasibility of extensive cross-fold validation and necessitating a single 80/20 stratified split.
    \item \textbf{Geographic Bias:} The dataset reflects the specific circulating serotypes and baseline hematocrit norms of South Asia, which limits direct transferability to populations in Latin America or Africa without re-calibration.
\end{enumerate}

\section{Conclusion and Future Work}
This study evaluated age-stratified machine learning models and subgroup-conditioned SHAP feature attributions for PLT-derived severe-dengue proxy prediction on a Bangladeshi clinical dataset ($n=972$). Stacking and Random Forest classifiers achieved predictive performance up to 90.6\% accuracy in pediatrics and 84.7\% accuracy in adults. Kendall's $\tau$-b rank correlation (0.900) confirmed high relative rank similarity overall, while depth-1 decision trees fitted to SHAP values identified age-dependent split thresholds (such as lymphocyte percentage splits of 44.50\% in children vs 24.99\% in adults).

Future work should focus on three areas: (1) prospective external validation of extracted feature splits on multi-center international cohorts, (2) expanding pediatric sample sizes to enable multi-fold cross-validation, and (3) integrating temporal laboratory parameters (day-of-illness tracking) alongside age stratification.

\begin{thebibliography}{99}
\bibitem{tanner2008} L. B. Tanner et al., "Decision Tree Algorithms Predict the Diagnosis and Outcome of Dengue Fever in the Early Phase of Illness," \textit{PLOS Negl. Trop. Dis.}, 2008.
\bibitem{potts2010} J. A. Potts et al., "Prediction of Dengue Disease Severity among Pediatric Thai Patients Using Early Clinical Laboratory Indicators," \textit{PLOS Negl. Trop. Dis.}, vol. 4, no. 8, p. e769, 2010.
\bibitem{rajapakse2014} S. Rajapakse et al., "Clinical features of dengue in adults and children: A comparative study," \textit{Trans. R. Soc. Trop. Med. Hyg.}, vol. 108, no. 1, pp. 24-29, 2014.
\bibitem{mendeley2026} A. Thakur and S. Nasrin, "Comprehensive Dengue Hematology and Clinical Dataset from Bangladesh for Machine Learning and Diagnostic Research," \textit{Mendeley Data}, V4, 2025, doi: 10.17632/jsbmtk8hty.4.
\bibitem{lundberg2017} S. M. Lundberg and S. I. Lee, "A unified approach to interpreting model predictions," in \textit{Proc. NeurIPS}, 2017.
\bibitem{niazi2026} M. Niazi and S. Momand, "Dengue Severity Level Prediction Using Hematology-Driven Machine Learning Models," \textit{Sci. J. Eng. Technol.}, vol. 3, no. 1, 2026.
\bibitem{breiman2001} L. Breiman, "Random Forests," \textit{Mach. Learn.}, vol. 45, no. 1, pp. 5-32, 2001.
\bibitem{pedregosa2011} F. Pedregosa et al., "Scikit-learn: Machine Learning in Python," \textit{J. Mach. Learn. Res.}, vol. 12, pp. 2825-2830, 2011.
\bibitem{chen2016} T. Chen and C. Guestrin, "XGBoost: A Scalable Tree Boosting System," in \textit{Proc. ACM SIGKDD}, 2016.
\bibitem{ke2017} G. Ke et al., "LightGBM: A Highly Efficient Gradient Boosting Decision Tree," in \textit{Proc. NeurIPS}, 2017.
\bibitem{wolpert1992} D. H. Wolpert, "Stacked generalization," \textit{Neural Netw.}, vol. 5, no. 2, pp. 241-259, 1992.
\bibitem{huang2020} S.-W. Huang et al., "Assessing the risk of dengue severity using demographic information and laboratory test results with machine learning," \textit{PLOS Negl. Trop. Dis.}, vol. 14, no. 12, p. e0008960, 2020.
\bibitem{chowdhury2022} S. U. Chowdhury et al., "Shapley-Additive-Explanations-Based Factor Analysis for Dengue Severity Prediction," \textit{J. Imaging}, vol. 8, no. 9, p. 229, 2022.
\bibitem{masum2025} A. K. M. Masum et al., "Enhancing Dengue Fever Diagnosis: A Machine Learning Framework with Stacking Ensemble and SHAP Explainability," \textit{IEEE QPAIN}, 2025.
\bibitem{tuhin2025} I. A. Tuhin et al., "An interpretable machine learning model for dengue detection with clinical hematological data," \textit{Healthc. Anal.}, 2025.
\bibitem{lu2025} B. Lu et al., "Assessing generalizability of a dengue classifier across multiple datasets," \textit{PLOS ONE}, vol. 20, no. 8, p. e0323886, 2025.
\bibitem{haque2025} M. E. Haque et al., "A novel interpretable and real-time dengue prediction framework using clinical blood parameters," \textit{Front. AI}, vol. 8, p. 1626699, 2025.
\bibitem{joly2024} N. A. Joly et al., "An Explainable AI Based Approach for Forecasting Dengue Diseases Using Patient's Phenotypic Traits," \textit{IEEE ICCIT}, 2024.
\bibitem{arrubla2023} W. Arrubla-Hoyos et al., "Comparison of classical machine learning and ensemble techniques in the context of dengue severity prediction," \textit{IEEE Col. Carib. Conf.}, 2023.
\bibitem{das2025} K. Das et al., "Optimized Feature-Driven Dengue Diagnosis Using Explainable Machine Learning Approaches," \textit{IEEE QPAIN}, 2025.
\bibitem{sangkaew2026} S. Sangkaew et al., "Early individualized risk prediction using clinical data for children during the febrile phase of dengue," \textit{PLOS Digit. Health}, vol. 5, p. e0001171, 2026.
\bibitem{madewell2024} Z. J. Madewell et al., "Machine learning for predicting severe dengue, Puerto Rico," \textit{medRxiv}, Nov. 2024.
\bibitem{halwala2025} R. Halwala et al., "Early prediction models for dengue in Sri Lanka: A comparative analysis," \textit{Int. J. Med. Inform.}, vol. 183, p. 105432, 2025.
\end{thebibliography}

\end{document}
